% !TeX root = ../../../../../main.tex
% chapters/sections/chapter3/subsections/policies/fiscal.tex

\subsection{財政政策}
\label{sec:chapter3_policies_fiscal}

\subsubsection{自国の財政政策}
\label{sec:chapter3_policies_fiscal_home}
政府は均衡予算を達成し税収のすべてを家計への一括移転\(t_t^H\)として還元する。
所得税率\(\tau_t^H\)はその定常状態の値\(\tau_{ss}^H\)の対数周囲での外生的な確率過程に従う。
\begin{equation}
\label{form:chapter3_policies_fiscal_home_tau_eq}
    \log(\tau_t^H)=(1-\rho_{\tau}^H)\log(\tau_{ss}^H)+\rho_{\tau}^H\log(\tau_{t-1}^H)+\varepsilon_t^{\tau,H}
\end{equation}
政府の予算制約式は以下である。
\begin{equation}
\label{form:chapter3_policies_fiscal_home_budget_eq}
    Nt_t^H=\sum_{h\in H}\tau_t^Hp_t^hy_t^h
\end{equation}
ここで付録A.4（定理4）において証明されるように総名目所得は以下の関係式で表すことができる。
\begin{equation}
\label{form:chapter3_policies_fiscal_nominal_income_identity_eq}
    \sum_{h\in H}p_t^hy_t^h=p_t^HY_t^H
\end{equation}
これらより集計レベルの予算制約式は以下のように決定される。
\begin{equation}
\label{form:chapter3_policies_fiscal_gov_budget_aggregate_eq}
    Nt_t^H=\tau_t^Hp_t^HY_t^H
\end{equation}
これより一人当たり移転額は以下のように決定される。
\begin{equation}
\label{form:chapter3_policies_fiscal_transfer_home_eq}
    t_t^H=\tau_t^H\frac{p_t^HY_t^H}{N}
\end{equation}

\subsubsection{外国の財政政策}
\label{sec:chapter3_policies_fiscal_foreign}
同様に外国政府の税率\(\tau_t^F\)は外生的な確率過程に従い、一人当たり移転額は以下のように決定される。
\begin{equation}
\label{form:chapter3_policies_fiscal_transfer_foreign_eq}
    t_t^{F*}=\tau_t^F\frac{p_t^{F*}Y_t^F}{M}
\end{equation}
\begin{equation}
\label{form:chapter3_policies_fiscal_tax_rule_foreign_eq}
    \log(\tau_t^F)=(1-\rho_{\tau}^F)\log(\tau_{ss}^{F*})+\rho_{\tau}^F\log(\tau_{t-1}^F)+\varepsilon_t^{\tau,F}
\end{equation}